% ============================================================
%  X-ray Scattering Simulations: A Practical Introduction
%  A tutorial guide for undergraduate and early PhD students
% ============================================================

\documentclass[nobib]{tufte-book}

% ---- Core packages ----
\usepackage[utf8]{inputenc}
\usepackage[T1]{fontenc}
\usepackage{microtype}
\usepackage{booktabs}
\usepackage{graphicx}
\usepackage{amsmath}
\usepackage{amssymb}

% ---- Code listings ----
\usepackage{listings}
\usepackage{xcolor}

% ---- Callout / info boxes ----
\usepackage{tcolorbox}
\tcbuselibrary{skins, breakable}

% ---- Hyperlinks ----
\usepackage{hyperref}
\hypersetup{
    colorlinks=true,
    linkcolor=black,
    urlcolor=blue!70!black,
    citecolor=black
}

% =============================================================
%  Color definitions
% =============================================================
\definecolor{codebackground}{RGB}{248,248,248}
\definecolor{codekeyword}{RGB}{0,0,180}
\definecolor{codestring}{RGB}{0,128,0}
\definecolor{codecomment}{RGB}{128,128,128}
\definecolor{calloutblue}{RGB}{220,235,252}
\definecolor{calloutblueborder}{RGB}{70,130,200}
\definecolor{calloutyellow}{RGB}{255,249,219}
\definecolor{calloutyellowborder}{RGB}{200,170,0}
\definecolor{calloutgreen}{RGB}{220,248,220}
\definecolor{calloutgreenborder}{RGB}{60,160,60}
\definecolor{calloutred}{RGB}{252,228,228}
\definecolor{calloutredborder}{RGB}{180,60,60}

% =============================================================
%  Code listing style
% =============================================================
\lstdefinestyle{pythonstyle}{
    language=Python,
    backgroundcolor=\color{codebackground},
    basicstyle=\ttfamily\small,
    keywordstyle=\color{codekeyword}\bfseries,
    stringstyle=\color{codestring},
    commentstyle=\color{codecomment}\itshape,
    numberstyle=\tiny\color{gray},
    numbers=left,
    stepnumber=1,
    numbersep=8pt,
    showstringspaces=false,
    breaklines=true,
    frame=single,
    framerule=0.4pt,
    rulecolor=\color{gray!40},
    xleftmargin=16pt,
    framexleftmargin=16pt,
    tabsize=4,
    captionpos=b,
}

\lstdefinestyle{bashstyle}{
    language=bash,
    backgroundcolor=\color{codebackground},
    basicstyle=\ttfamily\small,
    keywordstyle=\color{codekeyword}\bfseries,
    commentstyle=\color{codecomment}\itshape,
    numberstyle=\tiny\color{gray},
    numbers=none,
    showstringspaces=false,
    breaklines=true,
    frame=single,
    framerule=0.4pt,
    rulecolor=\color{gray!40},
    xleftmargin=16pt,
    framexleftmargin=16pt,
    tabsize=4,
}

\lstset{style=pythonstyle}

% =============================================================
%  Custom tcolorbox environments
% =============================================================

% Key concept box (blue)
\newtcolorbox{conceptbox}[1][]{
    enhanced,
    breakable,
    colback=calloutblue,
    colframe=calloutblueborder,
    fonttitle=\bfseries\small,
    title={Key Concept},
    #1
}

% Note / aside box (yellow)
\newtcolorbox{notebox}[1][]{
    enhanced,
    breakable,
    colback=calloutyellow,
    colframe=calloutyellowborder,
    fonttitle=\bfseries\small,
    title={Note},
    #1
}

% Try it yourself box (green)
\newtcolorbox{exercisebox}[1][]{
    enhanced,
    breakable,
    colback=calloutgreen,
    colframe=calloutgreenborder,
    fonttitle=\bfseries\small,
    title={Try It Yourself},
    #1
}

% Warning box (red)
\newtcolorbox{warningbox}[1][]{
    enhanced,
    breakable,
    colback=calloutred,
    colframe=calloutredborder,
    fonttitle=\bfseries\small,
    title={Watch Out},
    #1
}

% =============================================================
%  Title
% =============================================================
\title{Simulating X-ray Scattering:\\[6pt]
       A Practical Introduction\\[6pt]
       Using Python and GitHub}
\author{}
\date{}

% =============================================================
\begin{document}
% =============================================================

\maketitle
\tableofcontents

% ============================================================
%  PREFACE
% ============================================================
\chapter*{Preface}
\addcontentsline{toc}{chapter}{Preface}

These notes are intended for undergraduate students and researchers early in
their graduate careers who want to build practical computational skills while
exploring the physics of X-ray scattering from soft-matter systems.
The goal is deliberately modest: by the end of this guide you will be able
to write clean, well-documented Python code that simulates how X-rays scatter
from simple geometric objects---starting with a solid sphere and ending with a
multi-shell hollow vesicle that captures the essential physics of a lipid
bilayer membrane.
Along the way you will also learn how to manage software environments, track
your work with version control, and produce publication-quality figures.

No prior experience with scattering theory is assumed.
A working knowledge of algebra and basic calculus will be helpful, as will some
familiarity with programming in any language, though complete beginners to Python
are welcome.
Each section of this guide builds directly on the previous one, and it is
strongly recommended that you work through them in order.

The exercises at the end of each chapter are an essential part of the learning
process.
Full worked solutions are collected in Appendix~\ref{app:solutions}, but you
are strongly encouraged to attempt each problem before consulting them.

\bigskip
\noindent\textit{A note on software versions:}
This guide was written using Python~3.12 and \texttt{uv}~0.4.
The commands shown should work without modification on macOS, Linux, and
Windows (via PowerShell or WSL2), but minor differences may appear.
Where platform-specific notes are needed they are marked clearly.

% ============================================================
%  PART ONE
% ============================================================
\part{Setting Up Your Scientific Computing Environment}

% ------------------------------------------------------------
\chapter{Python and Package Management with \texttt{uv}}
\label{ch:python}
% ------------------------------------------------------------

\section{Why Python?}

Python has become the dominant language for scientific computing over the past
decade, and for good reason.
It is free and open-source, runs on every major operating system, and has an
enormous ecosystem of libraries purpose-built for numerical work, data
visualisation, and statistical analysis.
For our purposes the most important of these are:

\begin{itemize}
    \item \textbf{NumPy} --- fast numerical arrays and mathematical functions
    \item \textbf{SciPy} --- scientific algorithms including special functions
          and optimisation routines
    \item \textbf{Matplotlib} --- publication-quality plotting
    \item \textbf{JupyterLab} --- an interactive notebook interface that runs
          in your browser or directly inside VS Code
\end{itemize}

Python is also designed to be readable.
The goal of this guide is for you to understand every line of code you write,
not just copy it.
Keeping code readable---with clear variable names and short explanatory
comments---is a habit worth building from the very first session.

% ------------------------------------------------------------
\section{What is \texttt{uv} and Why Use It?}
% ------------------------------------------------------------

When you install Python, you get the language itself but not the scientific
libraries listed above.
Those must be installed separately, and managing them can become complicated
quickly: different projects may need different versions of the same library,
and installing everything into a single global environment is a reliable way
to create hard-to-diagnose conflicts.

\texttt{uv} is a modern Python package and project manager that solves this
cleanly.\sidenote{%
    \texttt{uv} is developed by Astral and is written in Rust, which makes it
    significantly faster than older tools like \texttt{pip} and
    \texttt{conda}.
    Full documentation is available at
    \url{https://docs.astral.sh/uv/}.}
It handles three things for you automatically:

\begin{enumerate}
    \item Installing and managing Python itself
    \item Creating isolated \emph{virtual environments} for each project
    \item Installing and tracking the exact versions of every package your
          project depends on
\end{enumerate}

The concept of a virtual environment deserves a moment of attention before
we continue.

\begin{conceptbox}[title={What is a Virtual Environment?}]
A virtual environment is an isolated copy of Python and its installed packages
that belongs to one specific project.
Think of it like a clean laboratory bench: everything on that bench is exactly
what that experiment needs, and nothing from a different experiment can
interfere with it.

When you work inside a virtual environment, installing or upgrading a package
affects only that project---not anything else on your computer.
This makes your work reproducible: a collaborator can recreate your exact
environment from a single configuration file and run your code without
modification.
\end{conceptbox}

% ------------------------------------------------------------
\section{Installing \texttt{uv}}
% ------------------------------------------------------------

Open a terminal.\sidenote{%
    \textbf{macOS:} the Terminal application is in
    \texttt{Applications\,/\,Utilities}.\\[4pt]
    \textbf{Windows:} use PowerShell (search the Start menu).
    Windows Subsystem for Linux (WSL2) is an excellent alternative that
    provides a full Linux environment inside Windows---see
    \url{https://learn.microsoft.com/en-us/windows/wsl/} for setup
    instructions.\\[4pt]
    \textbf{Linux:} any terminal emulator will work.}
Run the appropriate command for your operating system:

\medskip
\noindent\textbf{macOS and Linux:}
\begin{lstlisting}[style=bashstyle]
curl -LsSf https://astral.sh/uv/install.sh | sh
\end{lstlisting}

\medskip
\noindent\textbf{Windows (PowerShell):}
\begin{lstlisting}[style=bashstyle]
powershell -ExecutionPolicy ByPass -c "irm https://astral.sh/uv/install.ps1 | iex"
\end{lstlisting}

\medskip
After installation, close and reopen your terminal, then verify that
\texttt{uv} is available:

\begin{lstlisting}[style=bashstyle]
uv --version
\end{lstlisting}

You should see output like \texttt{uv 0.4.x}.
If you see a \texttt{command not found} error, consult the installation
guide at \url{https://docs.astral.sh/uv/getting-started/installation/}.

% ------------------------------------------------------------
\section{Creating Your First Project}
% ------------------------------------------------------------

\texttt{uv} organises work around \emph{projects}.
A project is a folder that contains your code alongside a configuration file
(\texttt{pyproject.toml}) that records which packages your project needs and
in which versions.

Navigate to wherever you store your work and create a new project:

\begin{lstlisting}[style=bashstyle]
cd ~/Documents        # or wherever you prefer to keep projects
uv init saxs-tutorial
cd saxs-tutorial
\end{lstlisting}

\texttt{uv} will create the following structure:

\begin{lstlisting}[style=bashstyle]
saxs-tutorial/
    pyproject.toml    # project configuration and dependencies
    README.md         # project description (edit this freely)
    hello.py          # a minimal starter script
\end{lstlisting}

\begin{notebox}[title={What is \texttt{pyproject.toml}?}]
\texttt{pyproject.toml} is the central configuration file for your project.
It records the project name, the Python version required, and the list of
packages your project depends on.
You will rarely need to edit it by hand---\texttt{uv} keeps it up to date
automatically as you add packages.
It is the file that allows anyone else (or a future you) to recreate your
environment exactly.
\end{notebox}

% ------------------------------------------------------------
\section{Installing Scientific Packages}
% ------------------------------------------------------------

Install the packages we will need throughout this guide:

\begin{lstlisting}[style=bashstyle]
uv add numpy scipy matplotlib jupyterlab
\end{lstlisting}

\texttt{uv} creates a virtual environment in a hidden \texttt{.venv} folder
inside your project, installs the requested packages into it, and updates
\texttt{pyproject.toml} automatically.
The first time you run this it may take a minute or two to download
everything.

Verify the installation by asking \texttt{uv} to run a quick Python check:

\begin{lstlisting}[style=bashstyle]
uv run python -c "import numpy, scipy, matplotlib; print('All packages OK')"
\end{lstlisting}

If you see \texttt{All packages OK}, your environment is ready.

\begin{exercisebox}
Open \texttt{pyproject.toml} in VS Code (we will install VS Code in
Chapter~\ref{ch:vscode}) and find the section that lists your installed
packages.
What version of NumPy was installed?
How would you find out what the latest available version is?
\end{exercisebox}

% ============================================================
\chapter{Version Control with Git and GitHub}
\label{ch:git}
% ============================================================

\section{Why Scientists Should Use Version Control}

Imagine writing a simulation, obtaining results you are happy with, then
making changes to improve the code---only to discover that the new version
no longer reproduces the original results, and you cannot remember exactly
what you changed.
This is an extremely common and genuinely frustrating experience.

Version control is the solution.
Git is a tool that records every change you make to your project files,
along with a short description of what changed and why.
At any point you can look back through the complete history of your work,
compare any two versions side by side, or restore an earlier state exactly.

GitHub is a website that stores a copy of your Git history in the cloud,
making it accessible from any computer and easy to share with collaborators
or a supervisor.\sidenote{%
    GitHub (\url{https://github.com}) is the most widely used hosting
    service in the scientific community, but alternatives exist---GitLab
    (\url{https://gitlab.com}) and Bitbucket (\url{https://bitbucket.org})
    offer similar functionality.}

\begin{conceptbox}[title={The Core Git Workflow}]
The everyday Git workflow involves three steps that you will repeat throughout
this guide and beyond:
\begin{enumerate}
    \item \textbf{Stage} the changes you want to save:\quad
          \texttt{git add}
    \item \textbf{Commit} with a short message describing what changed:\quad
          \texttt{git commit}
    \item \textbf{Push} your commit to GitHub:\quad
          \texttt{git push}
\end{enumerate}
Think of a commit as a labelled snapshot of your project at a specific moment.
A good commit message is short but informative.
\texttt{add sphere form factor function} is far more useful than
\texttt{update} or \texttt{changes} when you are looking back through
history six months later.
\end{conceptbox}

% ------------------------------------------------------------
\section{Installing Git}
% ------------------------------------------------------------

\noindent\textbf{macOS:}
Git is often pre-installed.
Check by running \texttt{git --version} in a terminal.
If it is absent, install the Xcode Command Line Tools:

\begin{lstlisting}[style=bashstyle]
xcode-select --install
\end{lstlisting}

Alternatively, download Git directly from \url{https://git-scm.com/downloads}.

\medskip
\noindent\textbf{Windows:}
Download the installer from \url{https://git-scm.com/downloads/win} and
accept the default options.
Git Bash, which is included with this installer, provides a convenient
terminal for the commands in this guide if you are not using WSL2.

\medskip
\noindent\textbf{Linux:}
\begin{lstlisting}[style=bashstyle]
# Debian or Ubuntu
sudo apt install git

# Fedora
sudo dnf install git
\end{lstlisting}

Confirm the installation:
\begin{lstlisting}[style=bashstyle]
git --version
\end{lstlisting}

% ------------------------------------------------------------
\section{Configuring Git}
% ------------------------------------------------------------

Before using Git for the first time, tell it your name and email address.
These are attached to every commit and will be visible in your project
history:

\begin{lstlisting}[style=bashstyle]
git config --global user.name  "Your Name"
git config --global user.email "your.email@example.com"
\end{lstlisting}

Use the same email address you will register with GitHub.

% ------------------------------------------------------------
\section{Creating a GitHub Account}
% ------------------------------------------------------------

Go to \url{https://github.com} and create a free account if you do not
already have one.
Choose a professional username---it will be visible to collaborators and
can appear on a CV or portfolio.

% ------------------------------------------------------------
\section{Setting Up SSH Authentication}
% ------------------------------------------------------------

GitHub needs a way to verify that it is really you when you push code to
your account.
The recommended method is SSH key authentication, which is more secure than
a password and, once configured, requires no further action from you.

\medskip
\noindent\textbf{Step 1: Generate an SSH key pair.}

\begin{lstlisting}[style=bashstyle]
ssh-keygen -t ed25519 -C "your.email@example.com"
\end{lstlisting}

When prompted for a file location, press Enter to accept the default
(\texttt{\textasciitilde/.ssh/id\_ed25519}).
When prompted for a passphrase, you may press Enter twice to leave it empty,
or set a passphrase for additional security.\sidenote{%
    On macOS, a passphrase can be stored in the system Keychain so you are
    not prompted repeatedly.}

\begin{conceptbox}[title={Public Keys and Private Keys}]
SSH authentication works by generating two mathematically linked files: a
\emph{private key} stored only on your computer (never shared) and a
\emph{public key} shared with GitHub.
When you connect to GitHub, your computer proves it holds the private key
without ever revealing it---similar to a lock-and-key system where the lock
is public but the key is yours alone.

The private key is stored at \texttt{\textasciitilde/.ssh/id\_ed25519}.
\textbf{Never share or upload your private key file under any circumstances.}
\end{conceptbox}

\medskip
\noindent\textbf{Step 2: Copy your public key.}

\begin{lstlisting}[style=bashstyle]
# macOS
cat ~/.ssh/id_ed25519.pub | pbcopy

# Linux (copy the printed output manually)
cat ~/.ssh/id_ed25519.pub

# Windows PowerShell
Get-Content ~/.ssh/id_ed25519.pub | Set-Clipboard
\end{lstlisting}

\medskip
\noindent\textbf{Step 3: Add the public key to GitHub.}

\begin{enumerate}
    \item Log in to \url{https://github.com}
    \item Click your profile picture (top right) and choose \textbf{Settings}
    \item In the left sidebar, click \textbf{SSH and GPG keys}
    \item Click \textbf{New SSH key}
    \item Give it a descriptive title (e.g.\ \textit{MacBook 2025}) and
          paste your public key into the Key field
    \item Click \textbf{Add SSH key}
\end{enumerate}

\medskip
\noindent\textbf{Step 4: Test the connection.}

\begin{lstlisting}[style=bashstyle]
ssh -T git@github.com
\end{lstlisting}

A successful connection produces a message like:\\[4pt]
\texttt{Hi username! You've successfully authenticated, but GitHub does not
provide shell access.}

\medskip
Full documentation for this process is available at
\url{https://docs.github.com/en/authentication/connecting-to-github-with-ssh}.

% ------------------------------------------------------------
\section{Creating Your Repository and Connecting It Locally}
% ------------------------------------------------------------

\noindent\textbf{Step 1: Create a new repository on GitHub.}

\begin{enumerate}
    \item Click the \textbf{+} icon (top right) and choose
          \textbf{New repository}
    \item Name it \texttt{saxs-tutorial}
    \item Add a short description if you like
    \item Choose \textbf{Public} or \textbf{Private} according to your
          preference
    \item Check \textbf{Add a README file}
    \item Click \textbf{Create repository}
\end{enumerate}

\medskip
\noindent\textbf{Step 2: Link your local project folder to GitHub.}

Inside your \texttt{saxs-tutorial} folder, run:\sidenote{%
    Replace \texttt{yourusername} with your actual GitHub username.}

\begin{lstlisting}[style=bashstyle]
git init
git remote add origin git@github.com:yourusername/saxs-tutorial.git
git branch -M main
\end{lstlisting}

You are now ready to push code from your computer to GitHub.

% ============================================================
\chapter{VS Code and Jupyter Notebooks}
\label{ch:vscode}
% ============================================================

\section{Installing VS Code}

Visual Studio Code (VS Code) is a free, cross-platform code editor developed
by Microsoft.
It has first-class support for Python and Jupyter notebooks and is widely
used in both academic research and industry.
Download it from \url{https://code.visualstudio.com} and follow the
installer for your operating system.

% ------------------------------------------------------------
\section{Installing Extensions}
% ------------------------------------------------------------

VS Code's capabilities are extended through \emph{extensions}.
Click the Extensions icon in the left sidebar (or press
\texttt{Ctrl+Shift+X} on Windows/Linux, \texttt{Cmd+Shift+X} on macOS)
and install:

\begin{itemize}
    \item \textbf{Python} (by Microsoft) --- language support, linting, and
          environment management
    \item \textbf{Jupyter} (by Microsoft) --- run Jupyter notebooks directly
          inside VS Code
\end{itemize}

After installing the Python extension, open the Command Palette
(\texttt{Ctrl+Shift+P} / \texttt{Cmd+Shift+P}) and search for
\texttt{Python: Select Interpreter}.
Choose the interpreter inside your \texttt{.venv} folder---it will appear
as something like \texttt{.venv (saxs-tutorial)}.
This tells VS Code to use the packages you installed with \texttt{uv}.

% ------------------------------------------------------------
\section{Opening Your Project}
% ------------------------------------------------------------

Go to \textbf{File $\rightarrow$ Open Folder} and select your
\texttt{saxs-tutorial} folder.
The file explorer on the left shows the project contents.

% ------------------------------------------------------------
\section{Anatomy of a Jupyter Notebook}
% ------------------------------------------------------------

A Jupyter notebook is a document made up of \emph{cells}, of which there
are two types you will use regularly:

\begin{itemize}
    \item \textbf{Markdown cells} --- formatted text, equations written in
          \LaTeX{} notation, and narrative explanation
    \item \textbf{Code cells} --- Python code that can be run interactively,
          one cell at a time, with output appearing directly below
\end{itemize}

This combination makes notebooks particularly well-suited to scientific work:
you can describe what you are about to compute, run the code, and have
the result---including figures---appear in the same document.
A well-written notebook is a complete record of a calculation that a
collaborator (or your future self) can follow and reproduce.

% ------------------------------------------------------------
\section{Hello, Scientific Python}
% ------------------------------------------------------------

Create a new file called \texttt{chapter\_01\_hello.ipynb}.
In the VS Code file explorer, click the new-file icon, type the filename
with the \texttt{.ipynb} extension, and VS Code will open it automatically
as a notebook.

Click \textbf{+ Markdown} to add your first markdown cell and type:

\begin{lstlisting}[style=bashstyle]
# Chapter 1: Hello, Scientific Python

This notebook confirms that the environment is working correctly.
\end{lstlisting}

Add a code cell and type the traditional first program:

\begin{lstlisting}
print("Hello, world!")
\end{lstlisting}

Press \texttt{Shift+Enter} to run the cell.
The text \texttt{Hello, world!} should appear immediately below it.

\medskip
Now test the scientific packages by adding a second code cell:

\begin{lstlisting}
import numpy as np
import matplotlib.pyplot as plt

# 200 evenly spaced values from 0 to 2*pi
x = np.linspace(0, 2 * np.pi, 200)
y = np.sin(x)

fig, ax = plt.subplots()
ax.plot(x, y, color="steelblue", label=r"$\sin(x)$")
ax.set_xlabel(r"$x$ (radians)")
ax.set_ylabel(r"$\sin(x)$")
ax.legend(loc="upper right")
plt.show()
\end{lstlisting}

Run this cell.
A sine wave should appear inline below the cell.
If it does, your environment is fully working and you are ready to move on.

\begin{notebox}[title={Inline Figures in VS Code}]
When running a plain \texttt{.py} script, matplotlib typically opens a
figure in a separate window.
Inside a Jupyter notebook in VS Code, figures appear inline, directly below
the code cell that produced them.
This is one of the reasons notebooks are convenient for exploratory
scientific work---the code, the narrative, and the results all live
together in one place.
\end{notebox}

% ============================================================
\chapter{Your First Commit}
\label{ch:firstcommit}
% ============================================================

You have set up your environment, installed your packages, and written your
first notebook.
Before moving to the science, take a few minutes to save this work
permanently using Git.
This habit---committing at every meaningful checkpoint---is one of the most
valuable practices you can build as a computational scientist.

% ------------------------------------------------------------
\section{Staging, Committing, and Pushing}
% ------------------------------------------------------------

Open a terminal inside VS Code by pressing the backtick key while holding
\texttt{Ctrl} (or going to \textbf{Terminal $\rightarrow$ New Terminal}).
You should be in your \texttt{saxs-tutorial} directory automatically.

\medskip
\noindent\textbf{Step 1: Check what Git can see.}

\begin{lstlisting}[style=bashstyle]
git status
\end{lstlisting}

Git will list any new or modified files that have not yet been committed.
Your notebook and \texttt{pyproject.toml} should appear here.

\medskip
\noindent\textbf{Step 2: Stage the files you want to include.}

\begin{lstlisting}[style=bashstyle]
git add chapter_01_hello.ipynb pyproject.toml README.md
\end{lstlisting}

\medskip
\noindent\textbf{Step 3: Commit with a clear message.}

\begin{lstlisting}[style=bashstyle]
git commit -m "add hello world notebook and project configuration"
\end{lstlisting}

\medskip
\noindent\textbf{Step 4: Push to GitHub.}

\begin{lstlisting}[style=bashstyle]
git push -u origin main
\end{lstlisting}

The \texttt{-u origin main} flags only need to be included on the first push.
After that, \texttt{git push} alone is sufficient.

Navigate to your repository on GitHub and refresh the page---your files
should now be visible there.

% ------------------------------------------------------------
\section{Ignoring Files That Should Not Be Tracked}
% ------------------------------------------------------------

Some files in your project should never be committed to Git.
The \texttt{.venv} folder, for instance, can be several hundred megabytes
and can always be recreated from \texttt{pyproject.toml} using
\texttt{uv sync}.
Compiled Python files (\texttt{.pyc}) are also unnecessary to track.

Create a file called \texttt{.gitignore} in your project root:

\begin{lstlisting}[style=bashstyle]
.venv/
__pycache__/
*.pyc
.DS_Store       # macOS folder metadata
\end{lstlisting}

Add this file to Git and commit it:

\begin{lstlisting}[style=bashstyle]
git add .gitignore
git commit -m "add gitignore for Python project"
git push
\end{lstlisting}

\begin{notebox}[title={GitHub's Python Template}]
When creating a new repository on GitHub, there is an option to add a
\texttt{.gitignore} automatically using a language-specific template.
Choosing the \textbf{Python} template provides a comprehensive list of
files to ignore and is a convenient starting point for new projects.
\end{notebox}

% ------------------------------------------------------------
\section{A Note on Commit Frequency}
% ------------------------------------------------------------

There is no single correct answer to how often you should commit.
A useful rule of thumb: commit whenever you reach a state you would be
unhappy to lose.
In practice this might mean committing after writing a new function, after
a plot looks the way you want it, or at the end of each working session.

Smaller, more frequent commits with descriptive messages are almost always
preferable to large, infrequent commits labelled \texttt{stuff} or
\texttt{final}.

\begin{exercisebox}
\begin{enumerate}
    \item Add a new markdown cell to your notebook that describes in your
          own words what the sine wave plot is showing.
    \item Run the notebook to confirm everything still works.
    \item Stage and commit the updated notebook with an appropriate message.
    \item Push the commit to GitHub and confirm it appears on the website.
\end{enumerate}
\end{exercisebox}

% ============================================================
%  APPENDICES
% ============================================================
\appendix

% ------------------------------------------------------------
\chapter{Remote Computing via SSH}
\label{app:remote}
% ------------------------------------------------------------

\emph{This appendix is provided for reference and can safely be skipped
for now.
You will not need remote computing until later in the tutorial series.
Return here when your supervisor or instructor indicates that it is time.}

\section{What is a Remote Computer and Why Use One?}

Some calculations in later chapters will be computationally expensive.
Rather than waiting for your laptop to finish a long job, you can connect
to a more powerful remote machine---such as a departmental server or a
high-performance computing cluster---and run the code there.

The standard tool for this is \textbf{SSH} (Secure Shell), which lets you
open a terminal session on a remote machine over a network connection.

% ------------------------------------------------------------
\section{Connecting to a Remote Machine}
% ------------------------------------------------------------

\begin{lstlisting}[style=bashstyle]
ssh username@hostname.university.edu
\end{lstlisting}

Replace \texttt{username} with your account name on the remote system and
\texttt{hostname.university.edu} with the address your institution provides.

If you have already set up an SSH key pair (Chapter~\ref{ch:git}), you can
copy your public key to the remote machine to avoid entering a password
each time:

\begin{lstlisting}[style=bashstyle]
ssh-copy-id username@hostname.university.edu
\end{lstlisting}

% ------------------------------------------------------------
\section{Transferring Files}
% ------------------------------------------------------------

To copy a file \emph{to} the remote machine:

\begin{lstlisting}[style=bashstyle]
scp localfile.ipynb username@hostname.university.edu:~/destination/
\end{lstlisting}

To copy a file \emph{from} the remote machine to your local computer:

\begin{lstlisting}[style=bashstyle]
scp username@hostname.university.edu:~/remotefile.ipynb ./
\end{lstlisting}

For whole folders, add the \texttt{-r} (recursive) flag:

\begin{lstlisting}[style=bashstyle]
scp -r saxs-tutorial/ username@hostname.university.edu:~/
\end{lstlisting}

% ------------------------------------------------------------
\section{Running Jupyter Notebooks on a Remote Machine}
% ------------------------------------------------------------

You can run a Jupyter server on the remote machine and view it in your
local browser by forwarding the network port over SSH.

\medskip
\noindent\textbf{Step 1:} Connect with port forwarding enabled:

\begin{lstlisting}[style=bashstyle]
ssh -L 8888:localhost:8888 username@hostname.university.edu
\end{lstlisting}

\noindent\textbf{Step 2:} On the remote machine, start JupyterLab:

\begin{lstlisting}[style=bashstyle]
uv run jupyter lab --no-browser --port=8888
\end{lstlisting}

\noindent\textbf{Step 3:} Copy the URL that JupyterLab prints (it will
look like \texttt{http://localhost:8888/lab?token=...}) and paste it
into your local browser.

You now have a full JupyterLab interface running on the remote machine,
displayed in your local browser.

% ------------------------------------------------------------
\section{Keeping Sessions Alive with \texttt{tmux}}
% ------------------------------------------------------------

If your SSH connection drops while a calculation is running, that
calculation will be killed.
\texttt{tmux} is a terminal multiplexer that keeps your session alive on the
remote machine even if the connection is interrupted.

\begin{lstlisting}[style=bashstyle]
# Start a new named session
tmux new -s tutorial

# Detach from the session (it keeps running in the background)
# Press Ctrl+B, then D

# Reattach after reconnecting via SSH
tmux attach -t tutorial
\end{lstlisting}

A concise \texttt{tmux} reference is available at
\url{https://tmuxcheatsheet.com}.

% ------------------------------------------------------------
\chapter{Exercise Solutions}
\label{app:solutions}
% ------------------------------------------------------------

\emph{Solutions will be added here progressively as the tutorial guide
grows.
Attempt each exercise on your own before consulting this appendix.}

\section*{Chapter~1}

\subsection*{Exercise 1.1}
Open \texttt{pyproject.toml} in VS Code.
The \texttt{[project]} section contains a \texttt{dependencies} key that
lists each installed package with its version constraint.
To see what the latest available version of NumPy is, run:

\begin{lstlisting}[style=bashstyle]
uv add numpy          # uv will report if a newer version is available
pip index versions numpy   # lists all available versions
\end{lstlisting}

The \texttt{uv} documentation at \url{https://docs.astral.sh/uv/} describes
how to upgrade packages to newer versions.

\section*{Chapter~4}

\subsection*{Exercise 4.1}
After adding your markdown cell and running the notebook, the staging and
commit sequence is:

\begin{lstlisting}[style=bashstyle]
git add chapter_01_hello.ipynb
git commit -m "add description of sine wave plot to notebook"
git push
\end{lstlisting}

The commit message should be specific to what changed.
If you also modified other files, include them in the \texttt{git add}
command before committing.

% ============================================================
\end{document}
